%% الفصل الأول — مدخل إلى التدبر
%% TODO: أضف محتوى الفصل هنا

\section{ما هو التدبر؟}

التدبر لغةً هو التفكر في عواقب الأمور. واصطلاحًا: إعمال العقل والقلب في فهم معاني القرآن الكريم واستخراج هداياته.

\begin{tadaburbox}
قال الله تعالى: ﴿أَفَلَا يَتَدَبَّرُونَ الْقُرْآنَ أَمْ عَلَىٰ قُلُوبٍ أَقْفَالُهَا﴾ (محمد: 24)
\end{tadaburbox}

\section{أهمية التدبر}

التدبر هو المقصود الأصلي من إنزال القرآن، كما قال تعالى: ﴿كِتَابٌ أَنزَلْنَاهُ إِلَيْكَ مُبَارَكٌ لِّيَدَّبَّرُوا آيَاتِهِ﴾.
